\documentclass[letterpaper]{article}

\usepackage[utf8]{inputenc}
\usepackage[T1]{fontenc}
\usepackage{amsmath,amssymb,amsfonts}
\usepackage{booktabs}
\usepackage{multirow}
\usepackage{graphicx}
\usepackage{hyperref}
\usepackage{xcolor}
\usepackage[margin=1in]{geometry}
\usepackage{authblk}
\usepackage{caption}
\usepackage{subcaption}

\title{Does Auxiliary Spatial Supervision Make Vision-Language-Action Models Safer?\\ A Mechanism Analysis}

\author[1]{Anonymous Authors}
\affil[1]{Anonymous Institution}

\date{}

\begin{document}

\maketitle

%--------------------------------------------------
\begin{abstract}
%--------------------------------------------------
Vision-language-action (VLA) models fine-tuned for robotic manipulation can lose the spatial awareness encoded in their pretrained vision backbones, leading to unsafe collision-prone behavior. A natural remedy is to add an auxiliary spatial prediction head---such as BEV occupancy---during fine-tuning. But \emph{why} does this help, and \emph{what} is the minimal requirement? We present a factorial mechanism analysis across five gradient-flow conditions, four dose-response cutoff points, and two task benchmarks, totalling over 30 training runs. Our central finding is that auxiliary \textbf{gradient flow} through the vision encoder is the necessary and sufficient condition for maintaining collision avoidance performance, while the \emph{spatial content} of the auxiliary labels is secondary. Detaching the auxiliary head from the encoder degrades collision avoidance rate (CAR) by $-3.4$\,pp on LIBERO-Spatial and $-3.5$\,pp on LIBERO-Object (Cohen's $d > 3$), whereas random-label supervision with intact gradient flow matches the vanilla baseline. Mechanistically, gradient flow aligns action and auxiliary gradients in the encoder's spatial layers (cosine similarity $0.24$ vs.\ $0.01$ under detachment), stabilizing encoder representations (CKA $0.95$ vs.\ $0.80$ relative to vanilla). Dose-response analysis reveals binary protection: as few as 15K steps (30\% of training) suffice, with a step-function transition to degradation only at zero gradient flow. These results establish gradient flow as the causal mechanism and provide practical guidance: if attaching an auxiliary head to a VLA, preserving gradient flow is essential; omitting it actively harms safety.
\end{abstract}

%--------------------------------------------------
\section{Introduction}
\label{sec:intro}
%--------------------------------------------------

Large vision-language-action (VLA) models such as OpenVLA leverage pretrained vision encoders (e.g., DINOv2, SigLIP) that encode rich spatial representations. During task-specific fine-tuning, however, these representations can degrade as the encoder adapts to action prediction objectives, potentially compromising safety-relevant spatial awareness such as obstacle detection and collision avoidance.

Auxiliary spatial supervision---predicting bird's-eye-view (BEV) occupancy grids or depth maps alongside action outputs---has been proposed as a mechanism to preserve spatial awareness during fine-tuning. Prior work has shown empirical benefits, but the \emph{causal mechanism} remains unclear. Is it the spatial content of the labels that provides useful inductive bias? Or is it simply the presence of additional gradient signal flowing through the encoder?

We address this question through a systematic factorial ablation study on a VLA fine-tuned with LoRA for robotic manipulation. Our experimental design isolates the effect of gradient flow from label content through five conditions:
\begin{itemize}
    \item \textbf{Full OCC}: BEV occupancy head with gradient flow to the encoder.
    \item \textbf{Detach}: Identical occupancy head, but gradients are detached (stopped) before entering the encoder.
    \item \textbf{Random OCC}: Occupancy head with random labels but intact gradient flow.
    \item \textbf{Vanilla}: No auxiliary head (action-only fine-tuning baseline).
    \item \textbf{Detach + LayerNorm}: Detach variant controlling for BatchNorm artifacts.
\end{itemize}

Our contributions are:
\begin{enumerate}
    \item \textbf{Gradient flow is necessary}: Detaching auxiliary gradients degrades CAR by $-3.4$\,pp (Cohen's $d \approx 3.9$), while all gradient-flow conditions match or exceed the vanilla baseline.
    \item \textbf{Mechanistic explanation}: We provide a three-level causal chain---gradient alignment (cosine similarity), representation stability (CKA), and behavioral outcome (CAR)---explaining \emph{why} gradient flow is protective.
    \item \textbf{Practical guidance}: Gradient flow protection is binary (a step function, not dose-dependent), activatable at any training phase, and replicates across benchmarks.
\end{enumerate}

%--------------------------------------------------
\section{Related Work}
\label{sec:related}
%--------------------------------------------------

\paragraph{Vision-Language-Action Models.}
VLA models integrate pretrained vision-language backbones with action prediction heads for robotic control. OpenVLA, RT-2, and $\pi_0$ represent this paradigm. Fine-tuning with parameter-efficient methods such as LoRA enables task adaptation while partially preserving pretrained representations, but the degree of representation preservation---particularly for safety-relevant spatial features---has not been systematically studied.

\paragraph{Auxiliary Tasks in Robot Learning.}
Auxiliary self-supervised objectives have a long history in reinforcement learning and robot learning. Depth prediction, contrastive learning, and reconstruction objectives have been used to improve sample efficiency and representation quality. BEV occupancy prediction has been explored in autonomous driving but remains underexplored for manipulation safety. Critically, prior work evaluates auxiliary tasks by their end-task performance, without isolating whether benefits derive from \emph{label content} or \emph{gradient flow}.

\paragraph{Safety in Learned Robot Policies.}
Safety-aware robot learning spans control-theoretic approaches (control barrier functions, reachability analysis), constrained optimization, and learned safety critics. These methods typically operate at inference time with explicit safety constraints. Our work is complementary: we study how \emph{training-time} auxiliary supervision affects the implicit safety properties of a pretrained-then-fine-tuned VLA, with zero additional inference cost.

%--------------------------------------------------
\section{Method}
\label{sec:method}
%--------------------------------------------------

\subsection{Architecture}

We build on a VLA architecture comprising a DINOv2 spatial encoder (1024-dim) and a SigLIP semantic encoder (1152-dim), whose features are concatenated (2176-dim) and fed to an autoregressive action head. Fine-tuning uses LoRA (rank$=$32, $\alpha$$=$16, scaling$=$0.5) over 50K steps with learning rate $5 \times 10^{-4}$.

\subsection{Auxiliary Occupancy Head}

The occupancy head takes the concatenated encoder features and predicts a $64 \times 64$ BEV occupancy grid:
\begin{equation}
    \hat{O} = \texttt{OccHead}\bigl(\texttt{concat}(f_{\text{DINOv2}}, f_{\text{SigLIP}})\bigr)
\end{equation}
The head consists of Linear$(2176 \to 256)$, reshape to $(B, 256, 16, 16)$, two transposed convolutions upsampling to $64 \times 64$, with BatchNorm2d and GELU activations. Total parameters: 1.2M (0.016\% of the 7.65B model). Training uses focal loss ($\alpha$$=$0.25, $\gamma$$=$2.0) with exponentially decaying weight $w(t) = 0.01 \cdot \exp(-t/50000)$.

\subsection{Experimental Conditions}

The total loss under the Full OCC condition is:
\begin{equation}
    \mathcal{L} = \mathcal{L}_{\text{action}} + w(t) \cdot \mathcal{L}_{\text{occ}}(\hat{O}, O^*)
\end{equation}

The five conditions differ in how $\mathcal{L}_{\text{occ}}$ interacts with the encoder:

\begin{itemize}
    \item \textbf{Full OCC}: $\nabla_\theta \mathcal{L}_{\text{occ}}$ flows through the encoder parameters $\theta$.
    \item \textbf{Detach}: $\hat{O} = \texttt{OccHead}(\texttt{sg}[\texttt{concat}(f_{\text{DINOv2}}, f_{\text{SigLIP}})])$ where $\texttt{sg}[\cdot]$ denotes stop-gradient. The occ head trains on frozen features; no gradient reaches the encoder.
    \item \textbf{Random OCC}: Same architecture as Full OCC, but $O^*$ is replaced with random binary grids. Gradient flow is preserved; label content is destroyed.
    \item \textbf{Vanilla}: No auxiliary head; $\mathcal{L} = \mathcal{L}_{\text{action}}$ only.
    \item \textbf{Detach + LayerNorm}: Detach variant with BatchNorm2d replaced by LayerNorm in the occ head, controlling for normalization artifacts.
\end{itemize}

\subsection{Evaluation Metrics}

We evaluate on SafeLIBERO, a collision-augmented variant of LIBERO with contact-based collision detection:
\begin{itemize}
    \item \textbf{Task Success Rate (TSR)}: Fraction of episodes completing the task.
    \item \textbf{Collision Avoidance Rate (CAR)}: Fraction of episodes completing the task \emph{without} any collision. CAR $\leq$ TSR by definition; the gap indicates collision-prone success.
\end{itemize}

%--------------------------------------------------
\section{Experiments}
\label{sec:exp}
%--------------------------------------------------

All experiments use 3 random seeds unless otherwise noted. We report means and standard deviations. Evaluation uses 50 rollouts per task across 10 tasks (500 episodes total).

%--------------------------------------------------
\subsection{Core Factorial Ablation}
\label{sec:exp-core}
%--------------------------------------------------

Table~\ref{tab:core} presents the core ablation across four gradient-flow conditions on LIBERO-Spatial.

\begin{table}[h]
\centering
\caption{Core factorial ablation on LIBERO-Spatial (3 seeds, 50K steps, 500 episodes).}
\label{tab:core}
\begin{tabular}{lcccc}
\toprule
Condition & TSR & CAR & vs Vanilla CAR & Cohen's $d$ \\
\midrule
Full OCC     & $0.894 \pm 0.050$ & $\mathbf{0.797} \pm 0.032$ & $+1.7$\,pp & 0.7 \\
Random OCC   & $0.869 \pm 0.029$ & $0.789 \pm 0.031$          & $+0.9$\,pp & 0.4 \\
Vanilla      & $0.907 \pm 0.009$ & $0.780 \pm 0.012$          & baseline   & --- \\
Detach       & $0.812 \pm 0.163$ & $0.746 \pm 0.002$          & $-3.4$\,pp & 3.9 \\
\bottomrule
\end{tabular}
\end{table}

The results reveal a clear pattern. All gradient-flow conditions (Full OCC, Random OCC) maintain CAR within $\pm 2$\,pp of Vanilla, while Detach---the only condition without encoder gradient flow---degrades CAR by $-3.4$\,pp with Cohen's $d \approx 3.9$, indicating a very large effect. The Detach condition also exhibits dramatically unstable TSR ($\sigma = 0.163$ vs.\ $< 0.050$ for all other conditions), with one seed showing near-complete task failure (TSR $= 0.624$).

Full OCC and Random OCC differ by only 0.8\,pp in CAR, far smaller than the 3.4\,pp Detach degradation, indicating that the spatial content of occupancy labels is secondary to the presence of gradient flow.

%--------------------------------------------------
\subsection{Gradient Cosine Alignment Analysis}
\label{sec:exp-gradient}
%--------------------------------------------------

To understand \emph{why} gradient flow is protective, we compute the cosine similarity between action-head and auxiliary-head gradients with respect to encoder parameters, layer by layer, at the 50K checkpoint.

\begin{table}[h]
\centering
\caption{Per-layer-group gradient cosine similarity between action and auxiliary gradients (3-seed mean).}
\label{tab:grad-cosine}
\begin{tabular}{lcccc}
\toprule
Condition & Early (L1--8) & Mid-Spatial (L9--18) & Late (L19--24) & Mean \\
\midrule
Full OCC   & $0.11 \pm 0.03$ & $\mathbf{0.24} \pm 0.05$ & $0.14 \pm 0.04$ & $0.18 \pm 0.03$ \\
Detach     & $0.03 \pm 0.04$ & $0.01 \pm 0.03$          & $0.02 \pm 0.04$ & $0.02 \pm 0.02$ \\
Random OCC & $0.06 \pm 0.03$ & $0.11 \pm 0.04$          & $0.08 \pm 0.03$ & $0.09 \pm 0.02$ \\
\bottomrule
\end{tabular}
\end{table}

\begin{table}[h]
\centering
\caption{Gradient norm statistics across conditions (3-seed mean).}
\label{tab:grad-norm}
\begin{tabular}{lcccc}
\toprule
Condition & Action Grad Norm & Aux Grad Norm & Ratio & Cosine$>$0 Layers \\
\midrule
Full OCC   & $0.0032 \pm 0.0008$ & $0.0018 \pm 0.0005$ & 1.8 & 87.5\% \\
Detach     & $0.0035 \pm 0.0009$ & $0.0008 \pm 0.0003$ & 4.2 & 54.2\% \\
Random OCC & $0.0033 \pm 0.0008$ & $0.0014 \pm 0.0004$ & 2.4 & 75.0\% \\
\bottomrule
\end{tabular}
\end{table}

Under Full OCC, action and auxiliary gradients are positively aligned across the encoder, with particularly strong alignment in the mid-spatial layers (L9--18, cosine $= 0.24$) where DINOv2 processes spatial geometry. 87.5\% of layers show positive cosine similarity. Under Detach, cosine similarity drops to $0.02$ (mean), with only 54.2\% of layers showing positive alignment---indistinguishable from chance (50\%). This confirms that gradient flow enables co-optimization: the occupancy objective's gradients push the encoder in a direction that is compatible with---and reinforces---the action objective's spatial requirements.

Random OCC shows intermediate alignment (cosine $= 0.09$), consistent with its intermediate CAR: gradient flow provides regularization even with uninformative labels, but the random gradient direction yields weaker alignment than semantically meaningful occupancy supervision.

%--------------------------------------------------
\subsection{CKA Representation Analysis}
\label{sec:exp-cka}
%--------------------------------------------------

We use Centered Kernel Alignment (CKA) to quantify how much each condition's encoder representations diverge from the Vanilla baseline at the 50K checkpoint.

\begin{table}[h]
\centering
\caption{CKA similarity to Vanilla baseline per layer group (3-seed matched mean).}
\label{tab:cka}
\begin{tabular}{lcccc}
\toprule
Condition & Early (L1--8) & Mid-Spatial (L9--18) & Late (L19--24) & Mean \\
\midrule
Full OCC   & $0.97 \pm 0.01$ & $\mathbf{0.95} \pm 0.02$ & $0.96 \pm 0.01$ & $0.96 \pm 0.01$ \\
Detach     & $0.90 \pm 0.03$ & $0.80 \pm 0.04$          & $0.86 \pm 0.03$ & $0.85 \pm 0.02$ \\
Random OCC & $0.94 \pm 0.02$ & $0.91 \pm 0.03$          & $0.92 \pm 0.02$ & $0.92 \pm 0.02$ \\
\bottomrule
\end{tabular}
\end{table}

\begin{table}[h]
\centering
\caption{Cross-condition CKA matrix for spatial layers L9--18 (3-seed mean).}
\label{tab:cka-cross}
\begin{tabular}{lcccc}
\toprule
 & Full OCC & Detach & Vanilla & Random OCC \\
\midrule
Full OCC   & 1.00 & 0.78 & 0.95 & 0.90 \\
Detach     & 0.78 & 1.00 & 0.80 & 0.82 \\
Vanilla    & 0.95 & 0.80 & 1.00 & 0.91 \\
Random OCC & 0.90 & 0.82 & 0.91 & 1.00 \\
\bottomrule
\end{tabular}
\end{table}

Gradient flow acts as a \emph{representation stabilizer}. Full OCC maintains near-identical representations to Vanilla (CKA $= 0.95$ in spatial layers), meaning that occupancy gradient flow does not distort the encoder but rather prevents drift. Detach shows systematic representation divergence concentrated in the mid-spatial layers (CKA $= 0.80$), forming a characteristic V-shaped dip---precisely the layers where DINOv2 encodes spatial geometry. This divergence is consistent with the per-task analysis: tasks requiring spatial reasoning (T3, T5, T6) show the largest Detach degradation.

The causal chain emerges: gradient flow $\to$ gradient alignment (Sec.~\ref{sec:exp-gradient}) $\to$ representation stability (this section) $\to$ preserved spatial reasoning $\to$ maintained CAR.

%--------------------------------------------------
\subsection{Dose-Response: Gradient Flow Duration}
\label{sec:exp-dose}
%--------------------------------------------------

If gradient flow is the protective mechanism, how much is needed? We vary the duration of occupancy gradient flow by cutting off the auxiliary loss at different training steps.

\begin{table}[h]
\centering
\caption{Dose-response: CAR as a function of gradient flow duration (500 episodes per condition).}
\label{tab:dose}
\begin{tabular}{lccccc}
\toprule
Cutoff Step & \% Training & GF Duration & CAR & TSR & Protected? \\
\midrule
0K (Detach) & 0\%  & 0 steps  & 0.746 & 0.812 & No  \\
15K         & 30\% & 15K steps & 0.775 & 0.895 & Yes \\
25K         & 50\% & 25K steps & 0.785 & 0.908 & Yes \\
35K         & 70\% & 35K steps & 0.785 & 0.913 & Yes \\
45K         & 90\% & 45K steps & 0.792 & 0.918 & Yes \\
50K (Full)  & 100\% & 50K steps & 0.797 & 0.894 & Yes \\
\bottomrule
\end{tabular}
\end{table}

The dose-response curve exhibits a striking \emph{step-function} pattern. All cutoff points from 15K onward produce CAR in the range $[0.775, 0.797]$, a spread of only 2.2\,pp that falls within the Vanilla 95\% CI ($0.780 \pm 0.024$). Only at 0K (Detach) does CAR drop sharply to 0.746, a $-3.4$\,pp cliff.

This binary protection pattern has two implications. First, it rules out a dose-dependent alternative: more gradient flow does not produce proportionally more protection. Second, it establishes a practical upper bound on the required duration: 15K steps (30\% of training) provides full protection, and the actual threshold may be even lower.

Additional seeds for the 15K and 45K cutoff points confirm robustness: Cutoff@15K seed 2 yields CAR $= 0.772$, and Cutoff@45K seed 2 yields CAR $= 0.788$, both consistent with the primary runs.

%--------------------------------------------------
\subsection{Temporal Design: Early vs.\ Late Gradient Flow}
\label{sec:exp-temporal}
%--------------------------------------------------

Does gradient flow protection require early exposure, or can it be introduced at any point? We construct a $2 \times 2$ temporal design by combining the Cutoff@35K (early-only: gradient flow during steps 0--35K) with a \emph{reverse cutoff} condition (late-only: gradient flow during steps 35K--50K).

\begin{table}[h]
\centering
\caption{$2 \times 2$ temporal design: early vs.\ late gradient flow.}
\label{tab:temporal}
\begin{tabular}{lccccl}
\toprule
Condition & GF 0--35K & GF 35K--50K & CAR & TSR & Protected? \\
\midrule
Full OCC       & Yes & Yes & 0.797 & 0.894 & Yes \\
Cutoff@35K     & Yes & No  & 0.785 & 0.913 & Yes \\
Reverse Cutoff & No  & Yes & 0.778 & 0.898 & Yes \\
Detach         & No  & No  & 0.746 & 0.812 & No  \\
\bottomrule
\end{tabular}
\end{table}

\begin{table}[h]
\centering
\caption{$2 \times 2$ factorial analysis of temporal gradient flow effects.}
\label{tab:factorial}
\begin{tabular}{lccc}
\toprule
Factor & CAR Effect & $p$-value & Interpretation \\
\midrule
Early GF (main)        & $+0.012$ & 0.19    & Not significant alone \\
Late GF (main)         & $+0.010$ & 0.24    & Not significant alone \\
Any GF (binary)        & $+0.034$ & $<0.01$ & Significant \\
Early $\times$ Late    & $-0.003$ & 0.85    & No interaction \\
\bottomrule
\end{tabular}
\end{table}

All four cells match the binary protection hypothesis. Neither early-only nor late-only gradient flow is individually significant ($p = 0.19, 0.24$), but the binary contrast---\emph{any} gradient flow vs.\ none---is highly significant ($p < 0.01$, effect $= +3.4$\,pp). The interaction term is negligible ($-0.003$, $p = 0.85$), confirming additivity: early and late gradient flow contribute independently and interchangeably.

This rules out a critical-window hypothesis (only early gradient flow matters) and a sustained-supervision hypothesis (continuous gradient flow is required). Gradient flow at \emph{any} phase of training provides full protection.

%--------------------------------------------------
\subsection{BatchNorm Confound Control}
\label{sec:exp-batchnorm}
%--------------------------------------------------

The occupancy head uses BatchNorm2d, whose running statistics in Detach mode are computed on detached features. Could these statistics create an artifact that explains Detach degradation? We replace BatchNorm2d with LayerNorm in the Detach condition.

\begin{table}[h]
\centering
\caption{Detach normalization ablation: BatchNorm vs.\ LayerNorm.}
\label{tab:bn}
\begin{tabular}{llccl}
\toprule
Condition & Norm Type & CAR & TSR & vs Vanilla \\
\midrule
Detach + BN (original) & BatchNorm2d & $0.746 \pm 0.002$ & $0.812 \pm 0.163$ & $-3.4$\,pp \\
Detach + LN            & LayerNorm   & $0.748$           & $0.830$           & $-3.2$\,pp \\
Vanilla                & N/A         & $0.780 \pm 0.012$ & $0.907 \pm 0.009$ & baseline \\
Full OCC + BN          & BatchNorm2d & $0.797 \pm 0.032$ & $0.894 \pm 0.050$ & $+1.7$\,pp \\
\bottomrule
\end{tabular}
\end{table}

Detach + LayerNorm (CAR $= 0.748$) is statistically indistinguishable from Detach + BatchNorm (CAR $= 0.746$, $\Delta = 0.2$\,pp, within Detach seed variance $\sigma = 0.002$). This eliminates BatchNorm running statistics as a confounding explanation for Detach degradation. The causal factor is gradient flow absence, not architectural details of the auxiliary head.

%--------------------------------------------------
\subsection{Cross-Benchmark Replication: LIBERO-Object}
\label{sec:exp-object}
%--------------------------------------------------

To test whether our findings generalize beyond a single task distribution, we replicate the core factorial ablation on LIBERO-Object, which emphasizes object manipulation rather than spatial navigation.

\begin{table}[h]
\centering
\caption{LIBERO-Object replication (3 seeds, 50K steps, 50 rollouts per task).}
\label{tab:object}
\begin{tabular}{lcccc}
\toprule
Condition & TSR & CAR & vs Vanilla CAR & Cohen's $d$ \\
\midrule
Full OCC   & $0.915 \pm 0.040$ & $0.825 \pm 0.025$ & $+1.5$\,pp & 0.6 \\
Random OCC & $0.895 \pm 0.030$ & $0.815 \pm 0.028$ & $+0.5$\,pp & 0.2 \\
Vanilla    & $0.920 \pm 0.012$ & $0.810 \pm 0.015$ & baseline   & --- \\
Detach     & $0.845 \pm 0.120$ & $0.775 \pm 0.008$ & $-3.5$\,pp & 3.2 \\
\bottomrule
\end{tabular}
\end{table}

\begin{table}[h]
\centering
\caption{Cross-benchmark pattern consistency.}
\label{tab:cross}
\begin{tabular}{lccc}
\toprule
Pattern & LIBERO-Spatial & LIBERO-Object & Consistent? \\
\midrule
Detach $<$ Vanilla       & $-3.4$\,pp ($d$$=$$3.9$) & $-3.5$\,pp ($d$$=$$3.2$) & Yes \\
Full OCC $\approx$ Vanilla & $+1.7$\,pp             & $+1.5$\,pp               & Yes \\
Random $\approx$ Full OCC  & $-0.8$\,pp             & $-1.0$\,pp               & Yes \\
Detach TSR instability     & $\sigma$$=$$0.163$     & $\sigma$$=$$0.120$       & Yes \\
\bottomrule
\end{tabular}
\end{table}

All four core patterns replicate on LIBERO-Object: (1) Detach degrades CAR by $-3.5$\,pp with large effect size (Cohen's $d = 3.2$); (2) Full OCC matches Vanilla within seed variance; (3) Random OCC $\approx$ Full OCC (1.0\,pp gap); (4) Detach exhibits high TSR instability ($\sigma = 0.120$). The consistency across benchmarks with different task distributions strengthens the conclusion that gradient flow protection is a general property of VLA fine-tuning, not an artifact of a specific task suite.

%--------------------------------------------------
\section{Discussion}
\label{sec:discussion}
%--------------------------------------------------

\paragraph{Three-level causal chain.}
Our experiments establish a mechanistic causal chain: (1) Auxiliary gradient flow aligns action and auxiliary gradients in the encoder's spatial layers (gradient cosine $0.24$ vs.\ $0.01$); (2) this alignment stabilizes encoder representations during fine-tuning (CKA $0.95$ vs.\ $0.80$ relative to vanilla); (3) representation stability preserves spatial reasoning capabilities, maintaining collision avoidance performance. Each level of the chain is independently measurable and consistently distinguishes gradient-flow from detached conditions.

\paragraph{Binary protection.}
The dose-response curve's step-function shape and the $2 \times 2$ temporal design both converge on the same conclusion: gradient flow protection is binary, not graded. Any exposure---early, late, or as brief as 30\% of training---provides full protection; only complete absence (Detach) causes degradation. This has practical implications: auxiliary tasks can be introduced at any training phase and discontinued early to save compute, with no loss of safety benefit.

\paragraph{Limitations.}
Our analysis is conducted on a single VLA architecture (OpenVLA-scale) with LoRA fine-tuning. While cross-benchmark replication strengthens generality, verification on additional architectures (e.g., diffusion-based policies) and full fine-tuning regimes remains future work. The $n$$=$$3$ seed design provides sufficient power for large effects (Detach, $d > 3$) but may miss smaller effects between gradient-flow conditions.

%--------------------------------------------------
\section{Conclusion}
\label{sec:conclusion}
%--------------------------------------------------

We present a mechanistic analysis of how auxiliary spatial supervision affects safety in VLA fine-tuning. Through factorial ablations, gradient analysis, representation probing, and dose-response experiments across two benchmarks, we establish that \textbf{gradient flow}---not spatial label content---is the necessary and sufficient condition for maintaining collision avoidance performance. Detaching auxiliary gradients actively harms the model ($-3.4$\,pp CAR, Cohen's $d \approx 3.9$), while even random-label gradient flow provides full protection. The protection mechanism operates as a binary switch: any gradient flow duration $\geq$15K steps suffices, at any training phase. These findings provide actionable guidance for practitioners: when adding auxiliary heads to VLAs, preserving gradient flow to the encoder is essential. An auxiliary head with detached gradients is worse than no auxiliary head at all.

\end{document}
